\documentclass[12pt]{article}
\usepackage{amsmath,amssymb,amsfonts}
\usepackage[margin=1in]{geometry}

\begin{document}

\section*{Chapter 3, Problem 8}

\textbf{Problem.} Prove that if the matrix $A$ is idempotent and is not the unit matrix, then $A$ is singular.

\bigskip
\textbf{Solution.}

A matrix $A$ is idempotent if $A^2 = A$.

Suppose for contradiction that $A$ is invertible. Then multiplying both sides of $A^2 = A$ on the left by $A^{-1}$:
\[
A^{-1} A^2 = A^{-1} A \quad \Longrightarrow \quad A = I.
\]

But we are told $A \neq I$. This is a contradiction. Therefore $A$ must be singular (non-invertible).

\medskip
\textit{Alternative proof using determinants:} From $A^2 = A$, taking determinants:
\[
\det(A^2) = \det(A) \quad \Longrightarrow \quad (\det A)^2 = \det A \quad \Longrightarrow \quad \det A\,(\det A - 1) = 0.
\]
So either $\det A = 0$ or $\det A = 1$. If $\det A = 1$, then $A$ is invertible, and by the argument above, $A = I$. Since $A \neq I$, we must have $\det A = 0$, so $A$ is singular. \hfill $\blacksquare$

\end{document}
